% Lösungen Punkte und Strecken im Koordinatensystem · Darstellen und Lage lesen · Prüfungs-Fokus Abitur GK (zusammenbau v0.6)
\einheitenkopf{Einheit 1 · Punkte darstellen und Lage lesen}

\erg{1}{a) $P$ liegt in der $x_1x_2$-Ebene \quad b) $C(4 | 5 | 2)$; $A$ und $B$ siehe Lösungsgrafik \quad c) Spitzen $(-2 | -2 | 4)$ und $(5 | 3 | 3)$; der erste Mast steht hinter der $x_2x_3$-Ebene und links der $x_1x_3$-Ebene; siehe Lösungsgrafik.}

1b)

\begin{ksys3} \rpunkt{5}{3}{1}{A} \rpunkt{2}{1}{4}{B} \rpunkt{4}{5}{2}{C} \end{ksys3}


1c)

\begin{ksys3} \rpunkt*{-2}{-2}{0}{F_1} \rpunkt*{5}{3}{0}{F_2} \rgerade[0:1]{-2,-2,0}{0,0,4}{} \rgerade[0:1]{5,3,0}{0,0,3}{} \rgerade[0:1]{-2,-2,4}{7,5,-1}{} \end{ksys3}

\erg{2}{a) Jeder Punkt hat zwei Koordinaten $0$ oder $4$ und liegt damit auf einer Würfelkante; die vier Punkte der Reihe nach verbinden, siehe Lösungsgrafik. \quad b) Jeder Punkt hat zwei Koordinaten $0$ oder $3$ und liegt damit auf einer Würfelkante; die vier Punkte der Reihe nach verbinden, siehe Lösungsgrafik. \quad c) Jeder Punkt hat zwei Koordinaten $0$ oder $4$ und liegt damit auf einer Würfelkante; die vier Punkte der Reihe nach verbinden, siehe Lösungsgrafik. \quad d) $P$ liegt auf der Deckkante mit $x_1 = 5$, ein Drittel ihrer Länge von vorn; Dreieck siehe Lösungsgrafik. \quad e) $P$ liegt auf der senkrechten Kante über $(6 | 0 | 0)$, ein Viertel der Höhe; Dreieck siehe Lösungsgrafik. \quad f) $P$ liegt auf der Deckkante mit $x_1 = 3$, ein Drittel ihrer Länge von vorn; Dreieck siehe Lösungsgrafik.}

2a)

\begin{ksys3} \rquader{0,0,0}{4}{4}{4} \rpunkt*{4}{0}{1}{P} \rpunkt*{1}{4}{0}{Q} \rpunkt*{0}{4}{2}{R} \rpunkt*{3}{0}{4}{S} \rgerade[0:1]{4,0,1}{-3,4,-1}{} \rgerade[0:1]{1,4,0}{-1,0,2}{} \rgerade[0:1]{0,4,2}{3,-4,2}{} \rgerade[0:1]{3,0,4}{1,0,-3}{} \end{ksys3}


2b)

\begin{ksys3} \rquader{0,0,0}{3}{3}{3} \rpunkt*{3}{1}{0}{P} \rpunkt*{3}{3}{2}{Q} \rpunkt*{1}{3}{3}{R} \rpunkt*{0}{0}{2}{S} \rgerade[0:1]{3,1,0}{0,2,2}{} \rgerade[0:1]{3,3,2}{-2,0,1}{} \rgerade[0:1]{1,3,3}{-1,-3,-1}{} \rgerade[0:1]{0,0,2}{3,1,-2}{} \end{ksys3}


2c)

\begin{ksys3} \rquader{0,0,0}{4}{4}{4} \rpunkt*{0}{3}{0}{P} \rpunkt*{4}{4}{1}{Q} \rpunkt*{4}{1}{4}{R} \rpunkt*{1}{0}{4}{S} \rgerade[0:1]{0,3,0}{4,1,1}{} \rgerade[0:1]{4,4,1}{0,-3,3}{} \rgerade[0:1]{4,1,4}{-3,-1,0}{} \rgerade[0:1]{1,0,4}{-1,3,-4}{} \end{ksys3}


2d)

\begin{ksys3} \rquader{0,0,0}{5}{6}{3} \rpunkt*{0}{0}{0}{O} \rpunkt*{5}{2}{3}{P} \rpunkt*{0}{6}{3}{R} \rgerade[0:1]{0,0,0}{5,2,3}{} \rgerade[0:1]{5,2,3}{-5,4,0}{} \rgerade[0:1]{0,6,3}{0,-6,-3}{} \end{ksys3}


2e)

\begin{ksys3} \rquader{0,0,0}{6}{3}{4} \rpunkt*{0}{0}{0}{O} \rpunkt*{6}{0}{1}{P} \rpunkt*{0}{3}{4}{T} \rgerade[0:1]{0,0,0}{6,0,1}{} \rgerade[0:1]{6,0,1}{-6,3,3}{} \rgerade[0:1]{0,3,4}{0,-3,-4}{} \end{ksys3}


2f)

\begin{ksys3} \rquader{0,0,0}{3}{6}{3} \rpunkt*{0}{0}{0}{O} \rpunkt*{3}{2}{3}{P} \rpunkt*{3}{6}{0}{T} \rgerade[0:1]{0,0,0}{3,2,3}{} \rgerade[0:1]{3,2,3}{0,4,-3}{} \rgerade[0:1]{3,6,0}{-3,-6,0}{} \end{ksys3}

\erg{3}{a) Die dritten Koordinaten $3$ und $2$ sind beide positiv: beide Punkte liegen oberhalb der $x_1x_2$-Ebene. \quad b) Alle Ecken haben $x_1 = 6$: parallel zur $x_2x_3$-Ebene. $B$ und $C$ unterscheiden sich nur im Vorzeichen von $x_2$, $A$ hat $x_2 = 0$: symmetrisch zur $x_1x_3$-Ebene. \quad c) Die dritten Koordinaten $-3$ und $2$ haben verschiedene Vorzeichen: $R$ liegt unterhalb, $T$ oberhalb der $x_1x_2$-Ebene. \quad d) Fehlende Ecken $(-5 | 3)$, $(-6 | 0)$, $(-5 | -3)$, $(0 | -4)$, $(5 | -3)$; siehe Lösungsgrafik. \quad e) Fehlende Ecken $(-3 | 2)$, $(-4 | 0)$, $(-3 | -2)$, $(0 | -3)$, $(3 | -2)$; siehe Lösungsgrafik. \quad f) Fehlende Ecken $(-2 | 5)$, $(-6 | 0)$, $(-2 | -5)$, $(0 | -6)$, $(2 | -5)$; siehe Lösungsgrafik.}

3d)

\begin{ksys}[xmin=-7,xmax=7,ymin=-7,ymax=7,klein] \punkt{6}{0}{} \punkt{5}{3}{} \punkt{0}{4}{} \punkt{-5}{3}{} \punkt{-6}{0}{} \punkt{-5}{-3}{} \punkt{0}{-4}{} \punkt{5}{-3}{} \end{ksys}


3e)

\begin{ksys}[xmin=-7,xmax=7,ymin=-7,ymax=7,klein] \punkt{4}{0}{} \punkt{3}{2}{} \punkt{0}{3}{} \punkt{-3}{2}{} \punkt{-4}{0}{} \punkt{-3}{-2}{} \punkt{0}{-3}{} \punkt{3}{-2}{} \end{ksys}


3f)

\begin{ksys}[xmin=-7,xmax=7,ymin=-7,ymax=7,klein] \punkt{6}{0}{} \punkt{2}{5}{} \punkt{0}{6}{} \punkt{-2}{5}{} \punkt{-6}{0}{} \punkt{-2}{-5}{} \punkt{0}{-6}{} \punkt{2}{-5}{} \end{ksys}

\erg{4}{a) $|SB| = |SD| = \sqrt{9 + 16} = 5$ cm (nicht $4$ cm); Spitzen bei $(-4 | 0)$ an $AD$, $(8 | 0)$ an $BC$ und $(0 | 8)$ an $CD$ – die Dreiecke an $BC$ und $CD$ sind bei $B$ bzw. $D$ rechtwinklig. \quad b) $|SB| = |SD| = \sqrt{16 + 9} = 5$ cm (nicht $3$ cm); Spitzen bei $(-3 | 0)$ an $AD$, $(9 | 0)$ an $BC$ und $(0 | 9)$ an $CD$ – die Dreiecke an $BC$ und $CD$ sind bei $B$ bzw. $D$ rechtwinklig. \quad c) $|SB| = |SD| = \sqrt{4 + 2{,}25} = 2{,}5$ cm (nicht $1{,}5$ cm); Spitzen bei $(-1{,}5 | 0)$ an $AD$, $(4{,}5 | 0)$ an $BC$ und $(0 | 4{,}5)$ an $CD$ – die Dreiecke an $BC$ und $CD$ sind bei $B$ bzw. $D$ rechtwinklig. \quad d) Die Grundfläche liegt parallel zur $x_1x_2$-Ebene, also zeigt $S'$ die Lage des Höhenfußpunkts: $C'(-1 | -1 | 0)$; $S'(4 | -1 | 0)$ liegt außerhalb des Dreiecks $A'B'C'$ – der Höhenfußpunkt liegt außerhalb der Grundfläche. \quad e) Die Grundfläche liegt parallel zur $x_1x_2$-Ebene, also zeigt $S'$ die Lage des Höhenfußpunkts: $C'(-1 | -2 | 0)$; $S'(1 | 1 | 0)$ liegt innerhalb des Dreiecks $A'B'C'$ – der Höhenfußpunkt liegt innerhalb der Grundfläche.}

4a)

\begin{ksys}[xmin=-5,xmax=9,ymin=-5,ymax=9,klein] \punkt{0}{0}{A} \punkt{3}{0}{B} \punkt{3}{3}{C} \punkt{0}{3}{D} \punkt{0}{-4}{S} \punkt{-4}{0}{S} \punkt{8}{0}{S} \punkt{0}{8}{S} \end{ksys}


4b)

\begin{ksys}[xmin=-4,xmax=10,ymin=-4,ymax=10,klein] \punkt{0}{0}{A} \punkt{4}{0}{B} \punkt{4}{4}{C} \punkt{0}{4}{D} \punkt{0}{-3}{S} \punkt{-3}{0}{S} \punkt{9}{0}{S} \punkt{0}{9}{S} \end{ksys}


4c)

\begin{ksys}[xmin=-2.5,xmax=5.5,ymin=-2.5,ymax=5.5,klein] \punkt{0}{0}{A} \punkt{2}{0}{B} \punkt{2}{2}{C} \punkt{0}{2}{D} \punkt{0}{-1.5}{S} \punkt{-1.5}{0}{S} \punkt{4.5}{0}{S} \punkt{0}{4.5}{S} \end{ksys}


4d)

\begin{ksys3} \rpunkt*{4}{1}{0}{A'} \rpunkt*{1}{5}{0}{B'} \rpunkt*{-1}{-1}{0}{C'} \rpunkt*{4}{-1}{0}{S'} \rgerade[0:1]{4,1,0}{-3,4,0}{} \rgerade[0:1]{1,5,0}{-2,-6,0}{} \rgerade[0:1]{-1,-1,0}{5,2,0}{} \end{ksys3}


4e)

\begin{ksys3} \rpunkt*{5}{-1}{0}{A'} \rpunkt*{1}{5}{0}{B'} \rpunkt*{-1}{-2}{0}{C'} \rpunkt*{1}{1}{0}{S'} \rgerade[0:1]{5,-1,0}{-4,6,0}{} \rgerade[0:1]{1,5,0}{-2,-7,0}{} \rgerade[0:1]{-1,-2,0}{6,1,0}{} \end{ksys3}

